\chapter{Limitaciones y amenazas a la validez}
\label{chap:limitaciones}

Este capítulo no es una disculpa genérica. Cada limitación se enuncia junto a la cifra concreta que
permite dimensionarla, porque una limitación que no puede cuantificarse no orienta al lector sobre
cuánto debe descontar de las conclusiones. La Tabla~\ref{tab:limitaciones} las reúne ordenadas por
severidad.

\begin{table}[H]
  \centering
  \footnotesize
  \caption{Limitaciones, evidencia que las cuantifica y severidad asignada.}
  \label{tab:limitaciones}
  % Tabla compuesta: cada fila enlaza una limitación con la cifra concreta que la
% cuantifica. Las cifras proceden de los artefactos citados en el capítulo:
% lo predictivo del Model Study 3 y lo económico del ganador del Portfolio Study.
\begin{tabularx}{\textwidth}{lXl}
\toprule
Limitación & Evidencia que la cuantifica & Severidad \\
\midrule
Multiplicidad de configuraciones &
Deflated Sharpe 0,682 sobre tres pasadas encadenadas más 1.728 carteras, por debajo del umbral
0,95 &
Alta \\
\addlinespace
Solapamiento de cohortes &
117 cohortes mensuales equivalen a $\approx 9$ observaciones independientes con horizonte de 12
meses &
Alta \\
\addlinespace
Potencia de la era reservada &
6 cohortes cerradas y $\approx 1$ observación independiente efectiva; 1,41 años de cartera y 1 de
2 años por encima del índice &
Alta \\
\addlinespace
Cartera elegida sobre la ventana de selección &
La ganadora es la mejor de 1.728: su IR de 0,844 en 2015--2024 es una cota superior optimista, no
una estimación insesgada &
Alta \\
\addlinespace
Dominio de un solo agente &
\textit{risk} aislado 0,1227 frente a 0,1090 del meta, que además le asigna más del 95\,\% del
peso; la neutralización retiene 86,62\,\% &
Media \\
\addlinespace
Mecanismo de la señal sin explicar &
El agente dominante se apoya en \texttt{gap\_21d} y \texttt{range\_63d} (78\,\% de las
observaciones), microestructura a semanas que ordena retornos a doce meses &
Media \\
\addlinespace
Alfa factorial no significativo en selección &
Alfa 0,27\,\% por periodo, $t=1{,}44$, $R^2=0{,}017$ &
Media \\
\addlinespace
Costes constantes &
5 pb de comisión y 10 pb de \textit{slippage}, sin impacto de mercado ni capacidad, sobre un
turnover del 324\,\% &
Media \\
\addlinespace
Universo restringido &
S\&P 500 estadounidense; 278 \textit{tickers} en 2003 frente a más de 500 en años recientes &
Media \\
\addlinespace
Versiones de catálogo heterogéneas &
Pasadas 1 y 2 bajo catálogo v6, pasada 3 bajo v7, que invierte un desempate por simplicidad &
Baja \\
\addlinespace
Salvaguarda \textit{a priori} en la curva de alfa &
No se activa en ninguna fila: la calibración usa horizonte (80\,\%) o era (20\,\%) &
Baja \\
\addlinespace
Factores no construibles &
Tamaño e inversión declarados no replicables en la regresión factorial &
Baja \\
\bottomrule
\end{tabularx}

  \caption*{La severidad es un juicio del autor sobre cuánto condiciona cada limitación las
  conclusiones, no una medida estadística.}
\end{table}

La principal limitación es la multiplicidad: se evaluaron 66 configuraciones y el Deflated Sharpe
es 0,930, inferior al umbral 0,95. Por ello el trabajo no afirma que la rentabilidad ajustada por
riesgo sobreviva concluyentemente a todas las pruebas de selección múltiple. Esta salvedad no anula
el Rank-IC, pero delimita la afirmación económica.

El horizonte de doce meses crea fuerte solapamiento entre cohortes: 117 observaciones mensuales
equivalen aproximadamente a nueve independientes. Además, la era reservada contiene sólo seis
cohortes cerradas y una observación independiente efectiva; su Rank-IC es indeterminado aunque la
confirmación económica sea favorable. El alfa factorial de la ventana de selección tampoco es
significativo por sí solo ($t=0,82$).

Existen límites de datos: universo restringido a Estados Unidos y al S\&P 500, cobertura desigual
de fundamentales, posibles restatements no observables en años antiguos y costes constantes que no
modelan impacto de mercado. Finalmente, el predominio de \textit{risk} obliga a interpretar con
prudencia la novedad de la señal; la neutralización conserva 84,35\%, lo que es evidencia parcial
contra una explicación puramente factorial, no una prueba definitiva.

Estas limitaciones conducen a una lectura disciplinada: la investigación demuestra una ordenación
predictiva robusta bajo un protocolo concreto y sugiere utilidad económica, pero requiere más
cohortes cerradas, universos adicionales y un catálogo más estrecho antes de extrapolarla.

\section{Amenazas a la validez interna}

El diseño point-in-time reduce de forma directa el \textit{lookahead bias}, pero depende de la
calidad de las fechas de filing y de la correspondencia entre métricas de Finnhub y periodos
contables. Los restatements históricos pueden ser invisibles en años tempranos: que una cifra esté
hoy disponible con una fecha de filing antigua no garantiza que sea idéntica a la cifra que conoció
el mercado ese día. La mitigación consiste en conservar periodo, filing y lag de ejecución, y en no
presentar la fuente como una cinta histórica perfecta.

La pertenencia histórica al índice limita el sesgo de supervivencia de composición, pero no elimina
la cobertura desigual de empresas retiradas ni la posible pérdida de fundamentales antiguos. El
estudio mide la fracción de filas utilizables y protege contra tickers reciclados; una réplica con
una fuente comercial point-in-time sería una prueba adicional de robustez. Esta limitación afecta
principalmente a la generalización, no a la trazabilidad del panel que sí se evaluó.

\section{Amenazas estadísticas y económicas}

Las 117 cohortes mensuales no son 117 ensayos independientes: el horizonte de doce meses hace que
dos cohortes contiguas compartan casi toda su ventana de retorno. Newey--West y bootstrap por bloques
reducen la sobreconfianza, pero no crean información nueva; por ello se informa también de unas
nueve observaciones efectivas. La multiplicidad añade una segunda cautela: las 66 configuraciones y
un Deflated Sharpe de 0,930 impiden una afirmación concluyente de Sharpe ajustado por selección.

Los costes son constantes, 5 pb de comisión y 10 pb de slippage. Esto demuestra que el resultado no
requiere coste cero, pero no modela impacto de mercado, bid--ask dinámico, capacidad ni fiscalidad.
Dado el turnover de selección, 359\% anualizado, una especificación dependiente de liquidez puede
modificar la cartera sin alterar el Rank-IC. El universo estadounidense de gran capitalización limita
además la extrapolación a pequeñas empresas, otros países o regímenes de liquidez diferentes.

La lectura correcta es específica: multiplicidad limita el Sharpe, solapamiento limita la precisión,
cobertura limita la generalización y cartera limita la transferencia de señal. Ninguna de estas
salvedades permite ignorar los placebos, la permutación o la estabilidad por eras; cada conclusión
debe mantenerse dentro del alcance de la evidencia que realmente la respalda.

\section{Limitaciones heredadas del proceso de desarrollo}

Tres limitaciones no proceden del diseño sino de decisiones tomadas durante el desarrollo. Merecen
figurar aquí porque afectan a la interpretación de las cifras y no serían visibles para un lector
que sólo viera el resultado final; su origen y la disyuntiva técnica que las motivó se documentan en
el Capítulo~\ref{chap:desarrollo-cartera}.

La primera es la salvaguarda de la curva de alfa: una recta fija entre $-10\,\%$ y $+10\,\%$ que se
activa en 4.384 de 20.545 filas del estudio, algo más de una de cada cinco. Toda cifra de cartera de
este trabajo está condicionada por ese mecanismo en una proporción apreciable de sus filas.

La segunda es que la calibración isotónica que precedió a esa salvaguarda se retiró sin conservarla
como opción del catálogo, de modo que este trabajo \emph{no puede} ofrecer una comparación empírica
entre ambas formulaciones; la afirmación de que la nueva es mejor descansa en un argumento sobre el
estimador, no en una medición pareada.

La tercera es que \texttt{minimum\_holding\_period} tomó el valor \texttt{none} por venir así en el
catálogo recomendado, no por haberse medido. El barrido diagnóstico del
Capítulo~\ref{chap:resultados} muestra que \texttt{half\_horizon} lo domina simultáneamente en
exceso, IR y rotación. No cambiarlo es metodológicamente correcto, porque las variables de cartera
no se optimizan; pero significa que la cartera evaluada no es la mejor que esta señal admite, y las
cifras económicas deben leerse como una cota inferior y no como el potencial del sistema.

\section{Reproducibilidad y congelación del estudio}

El estudio se ejecutó bajo la versión 5 del catálogo cerrado; el código actual está en la versión 6,
y por tanto no es reproducible bit a bit ni comparable con estudios de la versión 6. La razón de
congelar en lugar de reejecutar y el detalle de qué artefactos siguen siendo válidos se documentan
en la Sección~\ref{sec:reproducibilidad}. La consecuencia práctica para el lector es directa: un
tercero que ejecute hoy el código no obtendrá exactamente estos números.
